\section{Introducción}

En la actualidad, la propagación de enfermedades se encuentra estrechamente relacionada al flujo de información y, por ende, al de desinformación que hay en las redes sociales. Este fenómeno se hizo aún más evidente gracias a la pandemia de COVID-19 del 2020, siendo que la difusión de rumores, noticias falsas y teorías infundadas, junto con el pánico de las personas, las condicionan a comportarse de manera más escéptica sobre las medidas sanitarias tomadas por las autoridades. Este fenómeno, denominado ``infodemia'', que relaciona el contagio informativo y el biológico, no ha sido un suceso aislado. Giambiagi Ferrari et al. (2021) mostraron cómo el comportamiento de autoprotección de los seres humanos frente a una enfermedad evoluciona en función a lo que se diga sobre ella en su colectivo social, muy parecido a lo que ocurre con las opiniones o pensamientos. Sin embargo, a diferencia de los contagios por contacto, la información es capaz de circular de manera prácticamente instantánea gracias a las redes sociales. Esta disparidad de velocidades motiva el estudio de sistemas dinámicos de tipo rápido-lento para la modelización del comportamiento matemático de epidemias y el diseño de estrategias de comunicación.

En consecuencia de ello, resulta relevante demostrar que la evolución y trayectoria de una epidemia no depende únicamente de sus mecanismos biológicos, sino también de factores ajenos como la información. Acoplar un modelo SIRS con mecanismos informativos de distinta velocidad permite observar esta interacción y analizar su efecto sobre la población afectada (Bulai et al., 2024; Jia et al., 2019). Además, se necesita evaluar la separación de escalas para descubrir los métodos que mantengan estabilidad, positividad y precisión, obteniendo así simulaciones confiables bajo escenarios riesgosos.

Los modelos compartimentales provenientes del esquema de SIR se aplicaron inicialmente a enfermedades infecciosas, pero estudios posteriores extendieron su estructura a fenómenos de contagio no biológico. Por ejemplo, Jia et al. (2019) adaptaron el SIR a la propagación de rumores en redes sociales, distinguiendo dos canales de transmisión y analizando su derivación en $R_0 = 1$. De manera similar, Hernández Guillén y Martín del Rey (2018) lo extendieron a la propagación de malware, incorporando dispositivos ``portadores'' como vector de transmisión entre sistemas. Asimismo, Rehman et al. (2022) incrementaron la complejidad del enfoque incorporando memoria, recaída y reinfección en un modelo de compartimentos múltiples de malaria. Por su parte, Giambiagi Ferrari et al. (2021) acoplaron por primera vez un modelo epidemiológico con una dinámica continua de opiniones, mostrando que la estabilidad del equilibrio endémico dependía de la tolerancia al riesgo de la población. Sin embargo, ninguno de estos trabajos consideró que el comportamiento social y la enfermedad evolucionaran en escalas temporales distintas. Frente a ello, Bulai et al. (2024) resolvieron esta limitación proponiendo el modelo UMZ-SIRS, donde la desinformación evoluciona en una escala rápida y la enfermedad en una lenta, demostrando que la desinformación puede transformar una epidemia en extinción en una endemia persistente.

A pesar de los avances expuestos con anterioridad, y si bien Bulai et al. (2024) resolvieron la limitación conceptual de no distinguir escalas temporales al formular el modelo UMZ-SIRS, su análisis se concentra en el estudio cualitativo y geométrico del sistema mediante la Teoría Geométrica de Perturbaciones Singulares. Esto deja abierto un problema práctico: cuando dos procesos avanzan a velocidades tan distintas dentro de un mismo sistema de ecuaciones, este se vuelve numéricamente rígido, lo que dificulta seriamente su resolución computacional. Incluso trabajos recientes, como el de Li et al. (2024), que también acoplan información y epidemia, no analizan la estabilidad de los métodos numéricos cuando estas dos dinámicas operan en escalas de tiempo tan diferentes. Esto es relevante porque, si se resuelve este tipo de sistema con métodos tradicionales de paso fijo, como Runge-Kutta explícito, el resultado puede volverse inestable o incluso arrojar valores sin sentido físico, como poblaciones negativas. Evitar este problema exige pasos de tiempo extremadamente pequeños, lo cual encarece demasiado el cálculo computacional. Además, no existen suficientes estudios que comparen de forma sistemática métodos implícitos, como Euler hacia atrás, junto con solucionadores no lineales como Newton-Raphson, aplicados específicamente a sistemas epidemiológicos con esta separación de escalas.

Por ello, este proyecto propone modelar numéricamente el sistema SIRS rápido--lento acoplado a difusión de información y desinformación formulado por Bulai et al. (2024), con el propósito de analizar la estabilidad y persistencia de la epidemia frente a la separación de escalas temporales entre la dinámica informativa y la epidemiológica, y resolver la rigidez numérica que dicha separación introduce en el sistema de ecuaciones diferenciales mediante métodos explícitos e implícitos apropiados. De manera específica, este trabajo busca: 1) formular e implementar numéricamente el sistema UMZ-SIRS de seis ecuaciones diferenciales que representa el acoplamiento entre la dinámica rápida de información y desinformación y la dinámica lenta de la epidemia; 2) evaluar el comportamiento del modelo bajo un conjunto base de parámetros, comparando la respuesta de métodos numéricos explícitos e implícitos frente a la rigidez generada por la separación de escalas; 3) analizar la robustez del modelo mediante la variación de los parámetros de acoplamiento y de la tasa de contagio efectiva, cuantificando su efecto sobre la estabilidad numérica, la conservación de la población total y la aparición de rebrotes epidémicos; y 4) aplicar el modelo numérico para determinar las condiciones paramétricas bajo las cuales la difusión de desinformación transforma una epidemia en vías de extinción en una endemia persistente, contribuyendo al análisis de estrategias de comunicación en salud pública.
